\documentclass[10pt,a4paper]{amsart}
\usepackage[margin=19mm]{geometry}
\usepackage{amsmath,amssymb,mathtools}
\usepackage[scr=rsfs,cal=euler]{mathalfa}
\usepackage{xfrac}
\usepackage[unicode,hidelinks,bookmarks=false]{hyperref}
\newcommand{\ie}{i.e.\ }
\newcommand{\cf}{cf.\ }
\newcommand{\resp}{respectively\ }
\newcommand{\Subsection}{\subsection}
\providecommand{\nobreakdash}{\nobreak\hbox{-}}
\setlength{\emergencystretch}{2em}
\multlinegap0pt
\usepackage{latexsym,amscd}
\usepackage{enumerate}
\usepackage{graphicx}
\newcommand{\bsx}{\boldsymbol{x}}%
\newcommand{\bsm}{\boldsymbol{m}}%
\newcommand{\bsj}{\boldsymbol{j}}%
\newcommand{\bsy}{\boldsymbol{y}}%
\newcommand{\bsz}{\boldsymbol{z}}%
\newcommand{\bsk}{\boldsymbol{k}}%
\newcommand{\bsT}{\boldsymbol{T}}%
\newcommand{\bszero}{\boldsymbol{0}}%
\newcommand{\bsalpha}{\boldsymbol{\alpha}}%
\newcommand{\measure}[1]{\lambda_{s}(#1)}
\newcommand{\NN}{\mathbb N}
\newcommand{\ZZ}{\mathbb Z}
\newcommand{\RR}{\mathbb R}
\newcommand{\QQ}{\mathbb Q}
\newcommand{\PP}{\mathbb P}
\newcommand{\Log}{{\rm Log}}
\newcommand{\vol}{{\rm vol}}
\newcommand{\FF}{\mathbb{F}}
\renewcommand{\(}{\left (}
\renewcommand{\)}{\right )}
\newcommand{\digit}[2]{a_{#1}(#2)}
\newtheorem{theorem}{Theorem}[section]
\newtheorem{lemma}[theorem]{Lemma}
\newtheorem{remark}[theorem]{Remark}
\newtheorem{example}[theorem]{Example}
\newtheorem{prop}[theorem]{Proposition}
\newtheorem{coro}[theorem]{Corollary}
\newcommand{\notiz}{}
\theoremstyle{definition}
\newtheorem{prob}[theorem]{Open Problem}
\newtheorem{definition}[theorem]{Definition}
\newtheorem{conjecture}[theorem]{Conjecture}

\numberwithin{equation}{section}

\setcounter{section}{1}
\begin{document}
\title[Identifying van der Corput--Kronecker-type sequences as digital sequences]{Variants of the Littlewood conjecture, their connection to uniformly distributed sequences, and the exact order of the discrepancy of van der Corput--Kronecker-type sequences}
\author{Roswitha Hofer}
\date{Source: arXiv:2501.11362v2 (CC BY 4.0)}
\maketitle
\noindent\emph{Source and licence.} Variants of the Littlewood conjecture, their connection to uniformly distributed sequences, and the exact order of the discrepancy of van der Corput--Kronecker-type sequences. Version arXiv:2501.11362v2 (CC BY 4.0). This material is distributed under the Creative Commons Attribution 4.0 International licence, http://creativecommons.org/licenses/by/4.0/, as recorded in the arXiv OAI licence metadata for this exact version and shown on the abstract page https://arxiv.org/abs/2501.11362v2. Full rights notice: licenses/arxiv-2501.11362-CC-BY-4.0.txt.\par\smallskip
\noindent\textit{Editorial scope.} Counted unit: original Lemma 2.3 of the article (native environment \texttt{lemma}, source label \texttt{prop:1}, source0.tex lines 474--484, printed as Lemma 2.3 between Lemma 2.2 and Lemma 2.4 of the same section), the rank-structure lemma that recognises which upper-left submatrices of the Hankel matrix of a formal Laurent series are regular, delivered together with its complete original author proof as the single primary proof body at source lines 485--513. The proof environment follows the statement immediately and no detached proof block is involved. No proof marker is added and no mathematics is written, repaired or re-derived. Necessary complete context is retained uncounted: source lines 377--473, that is the whole of Section 2 from its own section heading, with the definition of the digital method and of digital \texttt{(t,s)}-sequences, the van der Corput sequence in base $p$ as a digital sequence, the discussion of continued fractions and best approximation denominators in $\FF_p((X^{-1}))$, the paragraph that explains why \texttt{prop:1} is needed, Lemma 2.2 with its complete proof, and the recognition of the van der Corput--Kronecker-type sequences as digital sequences. Nothing inside the delivered ranges is omitted except the seven forward cross-references described below. The publisher's own class is not present in the runtime, so the article is delivered in the AMS amsart wrapper of the pinned TeX Live runtime together with the author's own macro preamble, the author's own theorem declarations and the author's own equation numbering scheme, with the section counter restored editorially so that the delivered section prints as Section 2 and the theorem-like environments print as Definition 2.1, Lemma 2.2 and the counted Lemma 2.3, exactly as in the version PDF; the delivered ranges contain no numbered equation environment, so no equation number is affected. Seven cross-references of the delivered context point at results that lie outside the delivered unit, and they are delivered as link-only adaptations under a declared \texttt{reference\_only\_adaptation} receipt, each labelled with the number that the version PDF itself prints for that result: Theorem 1.8 (three occurrences), Theorem 1.10 (three occurrences) and Corollary 2.5. The article's own bibliography is not delivered as source text: the 12 works that the delivered ranges cite are re-typeset from the author's own inline \texttt{thebibliography} in the author's own order under the published numbers of that bibliography, so that the printed citation numbers coincide with the version PDF, namely 19, 23, 25, 27, 30, 33, 34, 37, 38, 39, 45 and 46. Printed slips kept literal, in particular the misspelling ``denomiators'' at source line 472 and the misspelling ``Laurentseries'' at source line 244, and the author's own redefinition of the math-shift control symbols \texttt{\textbackslash (} and \texttt{\textbackslash )} as \texttt{\textbackslash left(} and \texttt{\textbackslash right)}, which is carried into the wrapper verbatim. No figure, table, drawing or external asset occurs anywhere in the delivered ranges, so no artwork of any kind is redistributed and no bounded source comparison was needed.
\section{Identifying van der Corput--Kronecker-type sequences as digital sequences and a short proof of Theorem~\href{https://arxiv.org/abs/2501.11362v2}{1.8}}\label{sec:prerequ}

For some aspects of the class of van der Corput--Kronecker-type sequences, it is useful to interpret these sequences as digital $(t,s)$-sequences. In particular when investigating the discrepancy of a sequence the digital $(t,s)$-sequences concept, if applicable, is often the prefered method. We first define the digital method together with digital $(t,m,3)$-nets and digital $(t,2)$-sequences in the sense of Niederreiter (see e.g. \cite{niesiam}) as simple as possible to serve our purposes for Theorem~\href{https://arxiv.org/abs/2501.11362v2}{1.8} and Theorem~\href{https://arxiv.org/abs/2501.11362v2}{1.10}. We do not distinguish between $\FF_p$ and the set $\{0,1,\ldots,p-1\}$. 



\begin{definition}\label{def:t}
Let $p\in\PP$ and $m\in\NN$. Choose three $\NN\times m$-matrices $C_1,\,C_2,\,C_3$ over $\FF_p$, with $\NN$ indicating that they have infinitely many rows indexed by $1,2,3,\ldots$.  We construct a so-called digital pointset $(\bsx_n)_{0\leq n<p^m}$ of $p^m$ points as follows. Represent the nonnegative integer $n<p^m$ in base $p$, i.e.
$$ n = n_0+n_1p+\dots+n_{m-1}p^{m-1}\quad \mbox{with }0\leq n_j<p, $$
and set
$$ \vec n :=(n_0,\,\dots,\,n_{m-1})^T. $$
To generate the $i$th coordinate $x_n^{(i)}$ of $\bsx_n$, compute 
$$ C_i\cdot\vec n =:(y_1^{(i)},\,y_2^{(i)},\,\dots)^T\pmod{p}, $$
and set
$$ x_n^{(i)}:=\frac {y_1^{(i)}}p+\frac {y_2^{(i)}}{p^2}+\dots\,. $$
Let $t\in\NN_0$, $t<m$. We call the digital pointset $(\bsx_n)_{0\leq n<p^m}$ a \emph{digital $(t,m,3)$-net over $\FF_p$} if the following holds: 
%the parameter $\bsT$ is defined as follows. 
for all $d_1,d_2,d_3\in\NN_0$ with $d_1+d_2+d_3\leq m-t$ the $(d_1+d_2+d_3)\times m$-matrix consisting of the

upper $d_1$ rows of $C_1$ together with the

upper $d_2$ rows of $C_2$ together with the

upper $d_3$ rows of $C_3$, 

\noindent has full row rank $d_1+d_2+d_3$.


For a digital sequence $(\bsx_n)_{n\geq 0}$ in $ [0,1)^2$ we choose two $\NN\times\NN_0$-matrices $C_1,\,C_2$ over $\FF_p$  with $\NN_0$ indicating that they have infinitely many columns indexed by $0,1,2,3,\ldots$. To generate the $i$th coordinate $x_n^{(i)}$ of $\bsx_n$, represent the integer $n\geq 0$ in base $p$, i.e.
$$ n = n_0+n_1p+\dots+n_rp^r\quad \mbox{with }0\leq n_j<p, $$
and with $n_r\neq 0$, set
$$ \vec n :=(n_0,\,\dots,\,n_r,\,0,\,0,\,\dots)^T, $$
and 
$$ C_i\cdot\vec n =:(y_1^{(i)},\,y_2^{(i)},\,\dots)^T\pmod{p}. $$
Further
$$ x_n^{(i)}:=\frac {y_1^{(i)}}p+\frac {y_2^{(i)}}{p^2}+\dots\,. $$
Let $t\in\NN_0$. We call the digital sequence $(\bsx_n)_{n\ge 0}$ a \emph{digital $(t,2)$-sequence over $\FF_p$} if the following holds: 
%the parameter $\bsT$ is defined as follows. 
for every $m\in\NN, m>t$ and for all $d_1,d_2\in\NN_0$ with $d_1+d_2\leq m-t$ the $(d_1+d_2)\times m$-matrix consisting of the

left upper $d_1 \times m$-submatrix of $C_1$ together with the

left upper $d_2 \times m$-submatrix of $C_2$ 

\noindent has full row rank $d_1+d_2$.

\end{definition}

Usually the construction for digital sequences and digital nets is described for a general dimension $s$ and digital nets are based on $m\times m$ matrices instead of $\NN\times m$ matrices. For our proof of Theorem~\href{https://arxiv.org/abs/2501.11362v2}{1.10} we need more exact digital information. 
It is well-known that digital $(t,s)$-sequences are low-discrepancy sequences (cf. e.g. \cite[Theorem~4.17]{niesiam}). See for instance \cite{faure,Hofer,HoNie,nie,NieYeo,sobol,NieXing} for different constructions of $(t,s)$-sequences and \cite{Levints1,Levints2} for lower bounds of the discrepancy of the form $ND_N=\Omega(\log^s N)$ for many of them.  
 

%The main Theorem~\href{https://arxiv.org/abs/2501.11362v2}{1.10} of this manuscript ensures $ND_N=\Omega(\log^2 N)$ for the van der Corput--Kronecker-type sequences that are based on counterexamples of the $X$-adic Littlewood conjecture over a finite field $\FF_p$. %Before we do this, we need to clarify how these sequences are defined. 

The first step is to recognize the van der Corput--Kronecker-type sequences, which are based on counterexamples to the $X$-adic Littlewood conjecture over $\FF_p$, in Theorem~\href{https://arxiv.org/abs/2501.11362v2}{1.8} and in Theorem~\href{https://arxiv.org/abs/2501.11362v2}{1.10} as digital $(t,2)$-sequences over $\FF_p$. For this, it is important to explain these sequences within the concept of the digital method. We start with Lemma~\ref{ref:identifymatrices} which recognizes van der Corput--Kronecker-type sequences as digital sequences. 

\begin{lemma}\label{ref:identifymatrices}
Let $p\in\PP$ and $\theta=\sum_{i=j}^\infty a_iX^{-i}\in\FF_p((X^{-1}))$. Let $I$ denote the $\NN\times \NN_0$ unit matrix over $\FF_p$. Let $H(\theta)$ be the \emph{Hankel matrix} determined by the Laurent series $\theta$ defined as
$$H(\theta)=\begin{pmatrix} a_1&a_2&a_3&\ldots \\
	a_2&a_3&a_4&\ldots \\
	a_3&a_4&a_5&\ldots \\
	\vdots &\ddots&\ddots&\ddots\end{pmatrix}\in\FF_p^{\NN\times \NN_0},$$
 with the setting $a_1=a_2=\cdots=a_{j-1}=0$ in the case where $j>1$. Then the van der Corput--Kronecker-type sequence $((\varphi_p(n),{\left<  \theta n(X)\right>}|_p))_{n\geq 0}$ corresponds with the two-dimensional digital sequence generated by $I$ and $H(\theta)$. 
\end{lemma}
\begin{proof}
	Note that 
	$$\langle \theta n(X)\rangle=\langle \left(\sum_{i=1}^\infty a_iX^{-i}\right) \left(\sum_{r=0}^\infty n_rX^r\right)\rangle=\sum_{j=1}^\infty\left(\sum_{l=0}^\infty a_{j+l}n_l\right)X^{-j}. $$
	The rest is straightforward.
	
\end{proof}


%We remember, for a \emph{one-dimensional Kronecker type sequence} $(x_n)_{n\geq 0}$ first choose a Laurentseries $\theta=\sum_{i=w}^\infty a_iX^{-i}\in\FF_p((X^{-1}))$. We define a generating matrix %(We postpone basic information on Laurent series over $\FF_p$ to Section~\ref{sec:2} for the sake of readability of this section.) 
%As already intruduced in Section~\ref{sec:1} we write $n$ in base $p$, $n=n_0+n_1p+\cdots+n_rp^r$ associate the polynomial $n(X)\in\FF_p[X]$ as $n(X)=n_0+n_1X+\cdots+n_rX^r$. For the $n$th point compute the fractional part
%$\langle\theta n(X)\rangle$ and evaluate it by setting $X$ equal $p$. 
%It can be easily seen that the sequence $(\langle\theta n(X)\rangle|_p)_{n\geq 0}$ given by the coefficients of the Laurent series $\theta=\sum_{i=j}^\infty a_iX^{-j}$, can be interpreted as digital sequence with generating matrix $C$ given by $$C=\begin{pmatrix} c_{1,0}&c_{1,1}&c_{1,2}&\ldots \\
%c_{2,0}&c_{2,1}&c_{2,2}&\ldots  \\
%c_{3,0}&c_{3,1}&c_{3,2}&\ldots \\
%\vdots &\vdots&\vdots&\ddots\end{pmatrix}=\begin{pmatrix} a_1&a_2&a_3&\ldots \\
%a_2&a_3&a_4&\ldots \\
%a_3&a_4&a_5&\ldots \\
%\vdots &\ddots&\ddots&\ddots\end{pmatrix}=H(\theta)\in\FF_p^{\NN\times \NN_0}.$$
%If $j>1$ then we will set $a_1=a_2=\cdots=a_{j-1}=0$. The matrix $H(\theta)$ is called the \emph{Hankel matrix} determined by the Laurent series $\theta=\sum_{i=j}^\infty a_iX^{-i}$. % If $j>1$ then we will set $a_1=a_2=\cdots=a_{j-1}=0$. %(For detailed information on continued fractions and convergents in $\FF_p((X^{-1}))$ we refer the interested reader to the Appendix B of \cite{niesiam}.) 
%Straightforward one defines the $s$-dimensional Kronecker type sequence determined by $\theta_1,\ldots,\theta_s$ by just concatenating the one-dimensional Kronecker type sequences using $\theta_i$, i.e., $(\{n(X)\theta_1(X)\},\ldots,\{n(X)\theta_s(X)\})_{n\geq 0}$. 
%The \emph{van der Corput sequence in base $p$} can be defined using the $p$-adic radical inverse function or in the setting of digital sequences as the sequence that is obtained when using the $\NN\times \NN_0$ unit matrix $I$ over $\FF_p$ in Definition~\ref{def:t}. 


%Note that it had been an open problem whether there are low-discrepancy sequences constructed by $C_1=I$ and $C_2=H(\theta)$ or not. Finally, the recent counterexamples to the $X$-adic Littlewood conjecture by Adiceam et al. \cite{Adiceametal} gave a positive answer to this question. The fact that these counterexamples enable such low-discrepancy van der Corput--Kronecker-type sequences was first noticed by Levin \cite{levinLDhybrid} and later written down in more detail by Robertson \cite{Robertson}. 
The next step is to quantify the $t$-value for van der Corput--Kronecker-type sequences which are based on counterexample to the $X$-adic Littlewood conjecture. Previous work related the $t$-value of one-dimensional Kronecker-type sequences to the degrees of the coefficients in the simple continued fraction expansion of $\theta$ (see e.g. \cite{LarNie1993} and \cite{HoferdigKH}). Corollary~\href{https://arxiv.org/abs/2501.11362v2}{2.5} at the end of this section quantifies this $t$-value for van der Corput--Kronecker-type sequences, which are based on counterexample to the $X$-adic Littlewood conjecture, via a uniform bound for the degrees of the coefficients in the simple continued fraction expansions of $\langle X^r\theta\rangle,\,r\geq 0$. 
%Robertson raised the open problem for a lower bound for the star discrepancy of such a low-discrepancy hybrid sequence whose solution is the main task of this manuscript. Such a lower bound once again will support the important conjecture in the theory of uniform distribution, that for every sequence in $[0,1)^s$, $\limsup_{N\to\infty}ND_N^*/\log^sN >0$. 
%Let us first collect various aspects of a counterexample to the so-called $X$-adic Littlewood conjecture. For this purpose we recall and fix basic notation and information of simple continued fractions and convergents. 
%We write $\theta\in\FF_p((X^{-1}))$ as a series $\theta=\sum_{i=w}^\infty a_iX^{-i}\in\FF_p((X^{-1}))$ where $w\in\ZZ$ is minimal such that $a_w\neq 0$. The degree evaluation $\nu$ of $\theta$ is set $\nu(L)=-w$. The fractional part $\{\theta\}$ of $\theta$ is given by $\{\theta\}=\sum_{i=1}^\infty a_iX^{-i}$ and its polynomial part is set $\lfloor \theta\rfloor=\sum_{i=w}^0 a_iX^{-i}$. Completely analogously to $\RR$ 
For the sake of self-consistency, the basic properties of the simple continued fraction expansion of a formal Laurentseries over $\FF_p$ are summarized. We denote the simple continued fraction expansion of $\theta$ by $\theta=[A_0;A_1,A_2,\ldots]$ with polynomials $A_0=\theta-\langle\theta \rangle$ and $A_i\in\FF_p(X),\,i\geq 1$ of degrees $\geq 1$. % with $A_0=\lfloor\theta\rfloor$ and with polynomials $A_i\in\FF_p(X),\,i\geq 1$ of degrees $\geq 1$. Note that the expansion $\theta=[A_0;A_1,A_2,\ldots]$ is finite for rational $\theta$ and infinite else. 
For $h\geq 1$ the $h$th convergent $P_h/Q_h$ of $\theta$ is a rational function in $\FF_p(X)$ defined by $[A_0;A_1,\ldots,A_h]=:P_h/Q_h$, with $P_h,Q_h\in\FF_p[X]$ satisfying $\gcd(P_h,Q_h)=1$. The degree $\deg(Q_h)$ of $Q_h$ is often abbreviated to $d_h$ and satisfies $d_h=\sum_{i=1}^h\deg(A_i)$. Furthermore, $\deg(\theta-P_h/Q_h)=-d_h-d_{h+1}$ or equivalently $\deg(\langle Q_h \theta\rangle)=-d_{h+1}$ for every $h\geq 1$. For all $k\in\FF_p[X]$ with $d_h\leq \deg(k)<d_{h+1}$ we have $$\deg(\theta-b/k)\geq \deg(\theta-P_h/Q_h) \mbox{ for all $b\in\FF_p[X]$}$$ or equivalently $\deg(\langle k\theta\rangle)\geq \deg(\langle Q_h\theta\rangle)=-d_{h+1}$. In order to refer to the different magnitudes, as $A_i,\,d_h,\,Q_h$ etc. for $\langle X^r \theta\rangle$ instead of $\theta$ we add a superscript $(r)$. For example $[A^{(r)}_1,A^{(r)}_2,\ldots]$ denotes the simple continued fraction expansion of $\langle X^r \theta\rangle$ and $d^{(r)}_h=\sum_{i=1}^h\deg(A^{(r)}_i)$. (For concise information on continued fractions and convergents in $\FF_p((X^{-1}))$ we refer the interested reader to e.g. the Appendix B of \cite{niesiam}.)

Lemma~\ref{prop:1} gives an important property about the rank structure of the Hankel matrix $H(\theta)$ determined by a Laurent series $\theta\in\FF_p((X^{-1}))$. This property depends on the continued fraction of $\theta$ more exactly on the degrees of the best approximation denomiators. % =[A_0;A_1,A_2,\ldots]$ with $A_i\in\FF_{p}[X]$ for $i\geq 0$ and $\deg(A_i)\geq 1$ for $i\geq 1$. 


\begin{lemma}\label{prop:1}
	Let $\theta=\sum_{k=j}^\infty a_kX^{-k}$ be any formal Laurent series over $\FF_p$. Let $m\in\NN$. Then
	the upper left $m\times m$ submatrix of $H(\theta)$, 
	$$\begin{pmatrix}
		a_1&a_2&\cdots&a_m\\
		a_2&a_3&\cdots&a_{m+1}\\
		\vdots&\vdots&&\vdots\\
		a_{m}&a_{m+1}&\cdots&a_{m+m-1} 
	\end{pmatrix} $$
	over $\FF_p$ is regular if and only if $m=d_h=\deg(Q_h)$ for an $h\in\NN$. 
\end{lemma}

\begin{proof}
	We observe for a linear combination of the columns of the matrix that 
	$$b_0\begin{pmatrix}
		a_1\\
		a_2\\
		\vdots\\
		a_{m}
	\end{pmatrix} +b_1\begin{pmatrix}
		a_2\\
		a_3\\
		\vdots\\
		a_{m+1}
	\end{pmatrix}+\cdots + b_{m-1}\begin{pmatrix}
		a_m\\
		a_{m+1}\\
		\vdots\\
		a_{m+m-1}
	\end{pmatrix} =\begin{pmatrix}
		0\\
		0\\
		\vdots\\
		0
	\end{pmatrix}$$
	with $b_0,b_1,\ldots,b_{m-1}\in\FF_p$ if and only if 
	$$\deg(\langle(b_0+b_1X+\cdots+b_{m-1}X^{m-1})\theta\rangle)< -m.$$
	
	Let $m=d_h$. We know $\deg(\langle Q_{h-1}\theta\rangle)=-d_{h}$ and for all $P\in\FF_p[X]\setminus\{0\}$ with $\deg(P)<d_{h}$ we have $\deg(\langle P\theta\rangle)\geq \deg(\langle Q_{h-1}\theta\rangle)=-d_{h}$. From the observation above we derive the linear independence of the columns for such $m$. %Suppose that the columns are linearly dependent, then there exists a $P(X)=\sum_{r=0}^{d_h-1}p_rX^r$ not the zero polynomial such that $\nu(\{\thetaP\})<-d_h$, which is a contradiction. 
	Now if $d_{h-1}<m<d_h$ use $P=Q_{h-1}$ with $\deg P<m$ but $\nu(\{P\theta\})=-d_h$. Hence the columns are linearly depending. % to see the non-regularity of the matrix. 
\end{proof}

\begin{thebibliography}{99}
\bibitem[19]{faure} \textsc{H.~Faure}, \textit{Discr\'epance de suites associ\'ees \`a un syst\`eme de num\'eration (en dimension~$s$)}. {Acta Arith.} {{\bf41}} (1982), 337--351.
\bibitem[23]{HoferdigKH} \textsc{R.~Hofer}, \textit{Kronecker-Halton sequences in $\FF_p((X^{-1}))$}. {Finite Fields Appl.} {\bf50} (2018), 154--177.
\bibitem[25]{Hofer} \textsc{R.~Hofer}, \emph{A construction of low-discrepancy sequences involving finite-row digital $(t,s)$-sequences}. {{Math. Monatsh.}} {\bf171} (2013), 77--89.
\bibitem[27]{HoNie} \textsc{R.~Hofer and H.~Niederreiter}, \textit{A construction of $(t,s)$-sequences with finite-row generating matrices using global function fields}. {{Finite Fields Appl.}} {\bf21} (2013), 97--110.
\bibitem[30]{LarNie1993} \textsc{G.~Larcher, and H.~Niederreiter}, \textit{Kronecker-type sequences and non-Archimedean Diophantine approximations}. Acta Arith. {\bf63} (1993), no. 4, 379--396.
\bibitem[33]{Levints1} \textsc{M.B.~Levin}, \textit{On the lower bound of the discrepancy of $(t,s)$-sequences: I.} \emph{C. R. Math. Acad. Sci. Paris} \textbf{354} (2016), no. 6, 562--565.
\bibitem[34]{Levints2} \textsc{M.B.~Levin}, \textit{On the lower bound of the discrepancy of $(t,s)$-sequences: II,} {Online J. Anal. Comb. No.} {\bf12} (2017), 74 pp.
\bibitem[37]{nie} \textsc{H.~Niederreiter}, \textit{Low-discrepancy and low-dispersion sequences}. { J. Number Theory} {{\bf30}} (1988), 51--70.
\bibitem[38]{niesiam} \textsc{H.~Niederreiter}, \textit{Random Number Generation and Quasi-Monte Carlo Methods.} CBMS-NSF Regional Conference Series in Applied Mathematics, 63. SIAM, Philadelphia, 1992.
\bibitem[39]{NieYeo} \textsc{H.~Niederreiter and A.~S.~J.~Yeo}, \textit{Halton-type sequences from global function fields}. {{Sci. China Math.}} {\bf56} (2013), 1467--1476.
\bibitem[45]{sobol} \textsc{I.M.~Sobol'}, \textit{On the distribution of points in a cube and approximate evaluation of integrals}. {\u{Z}. Vy\v{c}isl. Mat. i Mat. Fiz.} {{\bf7}} (1967), 784--802.
\bibitem[46]{NieXing} \textsc{C.~Xing and H.~Niederreiter}, \textit{A construction of low-discrepancy sequences using global function fields}. {{ Acta Arith.}} {\bf73} (1995), 87--102.
\end{thebibliography}
\end{document}
